\documentclass[11pt,a4paper]{article}

% --- Essential Packages ---
\usepackage[top=2.5cm,bottom=2.5cm,left=2.5cm,right=2.5cm]{geometry}
\usepackage{amsmath,amssymb,amsfonts,mathtools,bm}
\usepackage{booktabs,tabularx,longtable,array,multirow}
\usepackage{xcolor}
\usepackage{listings}
\usepackage{tcolorbox}
\tcbuselibrary{skins,breakable}
\usepackage{enumitem}
\usepackage{fancyhdr}
\usepackage{microtype}
\usepackage{hyperref}

% --- Robust Unicode & Symbol Handling (pdflatex & xelatex/lualatex compatible) ---
\usepackage{iftex}
\ifPDFTeX
  \usepackage[utf8]{inputenc}
  \usepackage[T1]{fontenc}
  \usepackage{lmodern}
  \usepackage{textcomp}
  \usepackage{newunicodechar}
  \DeclareRobustCommand{\INR}{\textbf{Rs.\,}}
  \newunicodechar{₹}{\textbf{Rs.\,}}
  \newunicodechar{·}{\ensuremath{\cdot}}
  \newunicodechar{×}{\ensuremath{\times}}
  \newunicodechar{σ}{\ensuremath{\sigma}}
  \newunicodechar{–}{--}
  \newunicodechar{—}{---}
  \newunicodechar{←}{\ensuremath{\leftarrow}}
  \newunicodechar{→}{\ensuremath{\rightarrow}}
  \newunicodechar{↓}{\ensuremath{\downarrow}}
  \newunicodechar{−}{\ensuremath{-}}
  \newunicodechar{≥}{\ensuremath{\ge}}
  \newunicodechar{≤}{\ensuremath{\le}}
  \newunicodechar{┌}{\textsf{+}}
  \newunicodechar{┐}{\textsf{+}}
  \newunicodechar{└}{\textsf{+}}
  \newunicodechar{┘}{\textsf{+}}
  \newunicodechar{─}{\textsf{-}}
  \newunicodechar{│}{\textsf{|}}
  \newunicodechar{├}{\textsf{+}}
  \newunicodechar{┬}{\textsf{+}}
  \newunicodechar{►}{\ensuremath{\blacktriangleright}}
  \newunicodechar{◄}{\ensuremath{\blacktriangleleft}}
  \newunicodechar{ï}{\"{i}}
\else
  \usepackage{fontspec}
  \DeclareRobustCommand{\INR}{₹}
\fi

% --- Code Listing Configuration ---
\definecolor{codegreen}{rgb}{0,0.55,0.2}
\definecolor{codegray}{rgb}{0.45,0.45,0.45}
\definecolor{codepurple}{rgb}{0.55,0,0.75}
\definecolor{codeblue}{rgb}{0.05,0.2,0.65}
\definecolor{backcolour}{rgb}{0.97,0.97,0.99}

\lstset{
    backgroundcolor=\color{backcolour},   
    commentstyle=\color{codegreen}\itshape,
    keywordstyle=\color{codeblue}\bfseries,
    numberstyle=\tiny\color{codegray},
    stringstyle=\color{codepurple},
    basicstyle=\ttfamily\footnotesize,
    breakatwhitespace=false,         
    breaklines=true,                 
    captionpos=b,                    
    keepspaces=true,                 
    numbers=left,                    
    numbersep=6pt,                  
    showspaces=false,                
    showstringspaces=false,
    showtabs=false,                  
    tabsize=2,
    frame=single,
    rulecolor=\color{gray!30},
    literate=
      {┌}{{+}}1 {┐}{{+}}1 {└}{{+}}1 {┘}{{+}}1 {─}{{-}}1 {│}{{|}}1
      {├}{{+}}1 {┬}{{+}}1 {►}{{>}}1 {◄}{{<}}1 {↓}{{v}}1 {→}{{->}}2
      {←}{{<-}}2 {₹}{{Rs.}}2 {·}{{.}}1 {×}{{x}}1 {–}{{-}}1 {—}{{--}}1
      {≥}{{>=}}1 {≤}{{<=}}1
}

% --- Custom Column Types ---
\newcolumntype{Y}{>{\raggedright\arraybackslash}X}

% --- Hyperref Setup ---
\hypersetup{
    colorlinks=true,
    linkcolor=blue!70!black,
    citecolor=blue!70!black,
    filecolor=magenta,      
    urlcolor=blue!70!black,
}

% --- Header & Footer ---
\pagestyle{fancy}
\fancyhf{}
\rhead{\small\textit{Smart Evidence-based Triangulation for Uncovering Anomalies in MPLADS}}
\lhead{\small\textit{SETU --- Deep Technical Architecture}}
\cfoot{\thepage}
\renewcommand{\headrulewidth}{0.4pt}

% --- Title Metadata ---
\title{\Large\textbf{SETU --- Report 2: Data Pipeline \& Representation}\\[0.4em]\large\textit{Smart Evidence-based Triangulation for Uncovering Anomalies in MPLADS}}
\author{\textbf{SETU AI Core Team}\\Smart India Hackathon 2026 --- Ministry of Statistics and Programme Implementation (MoSPI)}
\date{September 2026}

\begin{document}
\maketitle
\tableofcontents
\vspace{1.5em}
\hrule
\vspace{1.5em}

\clearpage
\section*{REPORT 2 --- DATA PIPELINE \& DATA REPRESENTATION}
\addcontentsline{toc}{section}{REPORT 2 --- DATA PIPELINE \& DATA REPRESENTATION}

\section{What This Report Teaches}

Every transformation data undergoes from 12 raw relational CSVs to the final scaled feature matrix entering the ML models. Covers raw schemas, transaction aggregation mathematics, feature engineering (peer normalization), normalization (RobustScaler), score calibration, and complete data lineage.

\vspace{1em}\hrule\vspace{1em}

\section{Topic 6 --- Data Sources}

\subsection{The 12 Relational Tables}

\begin{table}[htbp]
\centering
\small
\begin{tabularx}{\textwidth}{c l c l Y}
\toprule
\textbf{\#} & \textbf{File} & \textbf{Rows} & \textbf{Grain} & \textbf{Role} \\
\midrule
01 & \texttt{01\_projects.csv} & 5,000 & 1 = 1 project & Central master --- work type, category, status, dates, coordinates \\
02 & \texttt{02\_financials.csv} & 5,000 & 1 = 1 project & Sanction amounts, SOR benchmarks, cost deviations, utilization ratios \\
03 & \texttt{03\_payments.csv} & \textbf{13,866} & 1 = 1 installment & Granular disbursements --- dates, amounts, MB flags, round-number flags \\
04 & \texttt{04\_progress.csv} & 5,000 & 1 = 1 inspection & Physical \%, financial \%, planned \%, geotag flags, MB verification status \\
05 & \texttt{05\_procurement.csv} & 5,000 & 1 = 1 project & Bidder counts, single-bid flags, tender duration, bid price similarity \\
06 & \texttt{06\_contracts.csv} & 5,000 & 1 = 1 project & Contract values, variation orders, extension counts, liquidated damages \\
07 & \texttt{07\_contractors.csv} & \textbf{400} & 1 = 1 contractor & Registration tier, blacklisting history, capacity limits, shared director flags \\
08 & \texttt{08\_agencies.csv} & \textbf{1,238} & 1 = 1 IDA & Agency type, staffing capacity, active project load \\
09 & \texttt{09\_geography.csv} & \textbf{619} & 1 = 1 district & Terrain type, literacy rate, poverty index, infrastructure gap score \\
10 & \texttt{10\_constituencies.csv} & \textbf{913} & 1 = 1 constituency & House type, MP allocation pool, total sanctioned funds \\
11 & \texttt{11\_documents.csv} & 5,000 & 1 = 1 project & AS/TS/MB/DPR/UC/CC flags, tampering flags, OCR confidence scores \\
12 & \texttt{12\_labels.csv} & 5,000 & 1 = 1 project & \textbf{QUARANTINED} --- is\_fraud, is\_hard\_negative, risk\_level, scenario\_type \\
\bottomrule
\end{tabularx}
\end{table}


\subsection{Key Relational Structure}

\begin{tcolorbox}[colback=gray!5!white,colframe=gray!60!black,arc=1mm,boxrule=0.6pt,left=6pt,right=6pt,top=6pt,bottom=6pt]
\begin{lstlisting}[basicstyle=\ttfamily\scriptsize,breaklines=true,columns=flexible]
01_projects (5,000)          ← Central anchor table
 ├── 1:1 ── 02_financials    ← Financial sanctions, cost deviations
 ├── 1:N ── 03_payments      ← MUST BE AGGREGATED (13,866 → 5,000 via aggregators.py)
 ├── 1:N ── 04_progress      ← MUST BE AGGREGATED (inspection rows → 1 per project)
 ├── 1:1 ── 05_procurement   ← Tender process metrics
 ├── 1:1 ── 06_contracts     ← Contract management data
 ├── N:1 ── 07_contractors   ← 400 contractors serve 5,000 projects
 ├── N:1 ── 08_agencies      ← 1,238 IDAs serve 5,000 projects
 ├── N:1 ── 09_geography     ← 619 districts → contextual enrichment
 ├── N:1 ── 10_constituencies← 913 constituencies → contextual enrichment
 ├── 1:1 ── 11_documents     ← Document audit trail
 └── 1:1 ── 12_labels        ← QUARANTINED — never joined into model inputs
\end{lstlisting}
\end{tcolorbox}


\subsection{Critical Schema Fields}

\textbf{From \texttt{02\_financials.csv}:}
\begin{itemize}[leftmargin=1.5em, itemsep=2pt, topsep=3pt]
\item \texttt{estimated\_cost} --- initial project estimate (baseline for cost deviation)
\item \texttt{cost\_per\_unit} --- normalized cost per unit of work (per metre for road, per seat for school)
\item \texttt{cost\_deviation} = (estimated\_cost − sanctioned\_amount) / sanctioned\_amount
\item \texttt{sor\_deviation} --- deviation from Ministry's Schedule of Rates benchmark
\item \texttt{fund\_utilization\_ratio} --- fraction of released funds actually spent
\end{itemize}

\textbf{From \texttt{03\_payments.csv} (transaction-grain, requires aggregation):}
\begin{itemize}[leftmargin=1.5em, itemsep=2pt, topsep=3pt]
\item \texttt{payment\_amount} --- installment size
\item \texttt{payment\_date} --- disbursement date (used for inter-payment interval computation)
\item \texttt{measurement\_book\_verified} --- boolean: was this payment MB-verified before release?
\item \texttt{is\_round\_number} --- boolean: does this payment end in 00,000?
\item \texttt{is\_final\_installment} --- boolean: is this the completion tranche?
\end{itemize}

\textbf{From \texttt{04\_progress.csv}:}
\begin{itemize}[leftmargin=1.5em, itemsep=2pt, topsep=3pt]
\item \texttt{physical\_progress} --- reported ground progress (\% complete)
\item \texttt{financial\_progress} --- disbursement progress (\% released)
\item \texttt{planned\_progress} --- scheduled progress
\item \texttt{geo\_tag\_available} --- GPS geotag present in this inspection?
\item \texttt{measurement\_book\_verified} --- MB verified at this inspection?
\end{itemize}

\vspace{1em}\hrule\vspace{1em}

\section{Topic 7 --- Raw Data $\rightarrow$ Usable Data}

\subsection{The Core Transformation Chain}

\[ X_\text{raw} \rightarrow X_\text{clean} \rightarrow X_\text{aggregated} \rightarrow X_\text{joined} \rightarrow X_\text{peer-normalized} \rightarrow X_\text{scaled} \]

\subsection{Step 1: Raw $\rightarrow$ Clean (\texttt{validators.py})}

Before any ML processing:
\begin{itemize}[leftmargin=1.5em, itemsep=2pt, topsep=3pt]
\item \textbf{Primary-Key Uniqueness}: Zero duplicate \texttt{project\_id} across all 12 tables
\item \textbf{Foreign-Key Integrity}: All 16 FK relationships checked --- zero orphan records
\item \textbf{Temporal Validity}: Dates parsed; invalid strings flagged
\item \textbf{Referential Loss}: 5,000 projects in $\rightarrow$ 5,000 rows out (zero lost)
\end{itemize}

\subsection{Step 2: Clean $\rightarrow$ Aggregated (\texttt{aggregators.py})}

This resolves the \textbf{many-to-one problem}.

\textbf{Payment Aggregation} (13,866 rows $\rightarrow$ 5,000 rows, 29 features):
\begin{tcolorbox}[colback=gray!5!white,colframe=gray!60!black,arc=1mm,boxrule=0.6pt,left=6pt,right=6pt,top=6pt,bottom=6pt]
\begin{lstlisting}[basicstyle=\ttfamily\scriptsize,breaklines=true,columns=flexible]
03_payments.csv (13,866 rows, one per installment)
    ↓  aggregate_payments()
project_payment_features (5,000 rows, 29 columns)
\end{lstlisting}
\end{tcolorbox}


\textbf{Progress Aggregation} (N inspection rows $\rightarrow$ 5,000 rows, 26 features):
\begin{tcolorbox}[colback=gray!5!white,colframe=gray!60!black,arc=1mm,boxrule=0.6pt,left=6pt,right=6pt,top=6pt,bottom=6pt]
\begin{lstlisting}[basicstyle=\ttfamily\scriptsize,breaklines=true,columns=flexible]
04_progress.csv (multiple inspection records per project)
    ↓  aggregate_progress()
project_progress_features (5,000 rows, 26 columns)
\end{lstlisting}
\end{tcolorbox}


\subsection{Step 3: Aggregated $\rightarrow$ Joined (\texttt{feature\_builder.py})}

Sequential controlled LEFT JOINs on \texttt{project\_id}:

\begin{lstlisting}[language=Python]
# Exact join sequence (feature_builder.py lines 197–210)
master = prepare_project_base(projects_df)           # 5,000 projects
master = master.merge(prefixed_financials, ...)       # + 28 financial features
master = master.merge(prefixed_payments, ...)         # + 29 payment features
master = master.merge(prefixed_progress, ...)         # + 26 progress features
master = master.merge(prefixed_procurement, ...)      # + 22 procurement features
master = master.merge(prefixed_contracts, ...)        # + 19 contract features
master = master.merge(contractor_context, ...)        # + 31 contractor features
master = master.merge(agency_context, ...)            # + 3 agency features
master = master.merge(geo_context, ...)               # + 8 geography features
master = master.merge(constituency_context, ...)      # + 1 constituency feature
master = master.merge(prefixed_documents, ...)        # + 50 document features
\end{lstlisting}


\textbf{Invariant checks enforced after join:}
\begin{lstlisting}[language=Python]
assert len(master) == 5000                     # No row duplication from joins
assert master["project_id"].nunique() == 5000  # No duplicate project_ids
assert len(dup_cols) == 0                      # No duplicate column names
for label_col in TARGET_LEAKAGE_COLUMNS:
    assert label_col not in master.columns     # ZERO label leakage
\end{lstlisting}


\subsection{Step 4: Identity Stripping (Anti-Overfitting)}

Before prefixing, entity identity strings are removed from model tensors:
\begin{lstlisting}[language=Python]
# feature_builder.py, extract_contractor_features(), lines 28–29
drop_cols = ["contractor_name", "home_district_id", "home_state_id"]
\end{lstlisting}

Also stripped: \texttt{agency\_name}, \texttt{mp\_name}, \texttt{district\_name}, \texttt{work\_name}, \texttt{title}.

This prevents the model from learning "Contractor X is always fraudulent" instead of learning structural behavioral patterns (which generalize to new contractors).

\vspace{1em}\hrule\vspace{1em}

\section{Topic 8 --- Feature Engineering}

\subsection{8.1 --- The Prefixing Convention}

Every feature in \texttt{project\_master\_features.csv} carries a domain prefix:

\begin{table}[htbp]
\centering
\small
\begin{tabularx}{\textwidth}{l l Y}
\toprule
\textbf{Prefix} & \textbf{Source} & \textbf{Example} \\
\midrule
\texttt{project\_\_} & 01\_projects.csv & \texttt{project\_\_work\_type} \\
\texttt{financial\_\_} & 02\_financials.csv & \texttt{financial\_\_estimated\_cost} \\
\texttt{payment\_\_} & 03\_payments.csv (aggregated) & \texttt{payment\_\_payment\_velocity\_mean} \\
\texttt{progress\_\_} & 04\_progress.csv (aggregated) & \texttt{progress\_\_financial\_minus\_physical\_progress} \\
\texttt{procurement\_\_} & 05\_procurement.csv & \texttt{procurement\_\_single\_bid\_flag} \\
\texttt{contract\_\_} & 06\_contracts.csv & \texttt{contract\_\_contract\_value\_change} \\
\texttt{contractor\_\_} & 07\_contractors.csv & \texttt{contractor\_\_blacklisted\_flag} \\
\texttt{agency\_\_} & 08\_agencies.csv & \texttt{agency\_\_active\_project\_count} \\
\texttt{geo\_\_} & 09\_geography.csv & \texttt{geo\_\_infrastructure\_gap\_index} \\
\texttt{constituency\_\_} & 10\_constituencies.csv & \texttt{constituency\_\_total\_sanctioned\_amount} \\
\texttt{document\_\_} & 11\_documents.csv & \texttt{document\_\_measurement\_book\_available} \\
\bottomrule
\end{tabularx}
\end{table}


\subsection{8.2 --- Payment Aggregation Mathematics}

For N installments per project, 29 behavioral statistics are computed:

\textbf{Volume \& Sizing:}
\[ \text{pay\_avg} = \frac{1}{N}\sum_{j=1}^{N} a_j, \qquad \text{pay\_std} = \sqrt{\frac{1}{N-1}\sum_{j=1}^{N}(a_j - \text{pay\_avg})^2} \]

\textbf{Payment Concentration (lump-sum fraud signal):}
\[ \text{pay\_conc\_max} = \frac{\max_j(a_j)}{\sum_j a_j} \]

\begin{tcolorbox}[colback=blue!4!white,colframe=blue!60!black,arc=1.5mm,boxrule=0.7pt,left=8pt,right=8pt,top=6pt,bottom=6pt]
A legitimate project has many moderate installments (concentration ~0.25–0.35). A ghost project often has one massive final installment equal to ~90\% of contract value (concentration ~0.90).\\
\end{tcolorbox}

\textbf{Payment Velocity (\INR{} per calendar day):}
\[ \text{pay\_velocity} = \frac{\sum_j a_j}{\max(1, t_\text{last} - t_\text{first})} \]

\begin{tcolorbox}[colback=blue!4!white,colframe=blue!60!black,arc=1.5mm,boxrule=0.7pt,left=8pt,right=8pt,top=6pt,bottom=6pt]
Ghost projects exhibit extremely high velocity --- large amounts released in a short window before reporting "complete."\\
\end{tcolorbox}

\textbf{Round Number Rate (structuring signal):}
\[ \text{round\_rate} = \frac{\text{count}(a_j \equiv 0 \pmod{100000})}{N} \]

\begin{tcolorbox}[colback=blue!4!white,colframe=blue!60!black,arc=1.5mm,boxrule=0.7pt,left=8pt,right=8pt,top=6pt,bottom=6pt]
Payments like \INR{}5,00,000 or \INR{}14,90,000 appear disproportionately in structuring scenarios (just below the \INR{}15L mandatory tender threshold).\\
\end{tcolorbox}

\textbf{Unverified Payment Rate:}
\[ \text{unverified\_rate} = \frac{\text{count}(\text{mb\_verified} = \text{False})}{N} \]

\textbf{Inter-Payment Interval Statistics:}
\[ \text{avg\_interval} = \frac{1}{N-1}\sum_{j=2}^{N}(t_j - t_{j-1}), \qquad \text{interval\_std} = \sigma(\{t_j - t_{j-1}\}_{j=2}^{N}) \]

\subsection{8.3 --- Progress Aggregation Mathematics}

\textbf{Physical-Financial Progress Gap (ghost-project signal):}
\[ \text{gap} = \text{financial\_progress} - \text{physical\_progress} \]

\begin{tcolorbox}[colback=blue!4!white,colframe=blue!60!black,arc=1.5mm,boxrule=0.7pt,left=8pt,right=8pt,top=6pt,bottom=6pt]
Example: 80\% funds disbursed, 5\% physical complete $\rightarrow$ gap = 0.75. The primary ghost project indicator.\\
\end{tcolorbox}

\textbf{Physical Progress Slope (time-aware, computed against calendar days):}
\[ \text{phys\_slope} = \frac{\text{physical\_progress}_\text{latest} - \text{physical\_progress}_\text{first}}{\max(1, t_\text{latest} - t_\text{first})} \]

\begin{tcolorbox}[colback=blue!4!white,colframe=blue!60!black,arc=1.5mm,boxrule=0.7pt,left=8pt,right=8pt,top=6pt,bottom=6pt]
Critical: slope is per \textbf{calendar day}, not row index. Row-index slope would be meaningless if inspections are unevenly spaced.\\
\end{tcolorbox}

\textbf{Gap Slope (growing divergence indicator):}
\[ \text{gap\_slope} = \frac{\text{gap}_\text{latest} - \text{gap}_\text{first}}{\text{span\_days}} \]

\begin{tcolorbox}[colback=blue!4!white,colframe=blue!60!black,arc=1.5mm,boxrule=0.7pt,left=8pt,right=8pt,top=6pt,bottom=6pt]
Positive increasing gap slope = spending money faster than completing physical work = growing divergence that signals fraud or misreporting.\\
\end{tcolorbox}

\subsection{8.4 --- Peer Deviation Engineering (the most critical step)}

These features transform absolute financial values into \textbf{peer-relative deviations}, allowing the Isolation Forest to detect fraud at any project scale.

\textbf{For metric $x_i$ in project $i$ (e.g., \texttt{financial\_\_estimated\_cost}):}

\textbf{Peer Difference:}
\[ \Delta_i = x_i - \text{Median}(X_{\text{peer}(i)}) \]

\textbf{Peer Ratio:}
\[ \rho_i = \frac{x_i + \varepsilon}{\text{Median}(X_{\text{peer}(i)}) + \varepsilon} \]

\textbf{Robust Z-Score (MAD-normalized):}
\[ Z_i^\text{robust} = \frac{x_i - \text{Median}(X_{\text{peer}(i)})}{1.4826 \times \text{MAD}(X_{\text{peer}(i)}) + \varepsilon} \]

Where:
\begin{itemize}[leftmargin=1.5em, itemsep=2pt, topsep=3pt]
\item $X_{\text{peer}(i)}$ = all projects in the same \texttt{work\_type} group
\item $\text{MAD} = \text{Median}(|X_\text{peer} - \text{Median}(X_\text{peer})|)$
\item $1.4826$ = standard normal consistency constant (makes MAD-based Z comparable to standard Z for Gaussian data)
\item $\varepsilon = 10^{-6}$ (numerical stability, from \texttt{config.py peer\_epsilon})
\item Falls back to global statistics if peer group count < 5
\end{itemize}

\textbf{Peer Groups:}
\begin{itemize}[leftmargin=1.5em, itemsep=2pt, topsep=3pt]
\item \textbf{Primary}: \texttt{project\_\_work\_type} (29 work types --- e.g., "Road Construction", "Drinking Water", "Sanitation")
\item \textbf{Secondary}: \texttt{project\_\_category} (8 broader categories --- fallback if work-type peer group is too small)
\item \textbf{Tertiary}: Global population (final fallback)
\end{itemize}

\textbf{All 11 engineered peer features:}

\begin{table}[htbp]
\centering
\small
\begin{tabularx}{\textwidth}{l l Y}
\toprule
\textbf{Feature} & \textbf{Formula} & \textbf{Business Signal} \\
\midrule
\texttt{peer\_\_cost\_diff\_work\_type\_median} & $x_i - \text{Median}$ & Raw cost excess vs peer \\
\texttt{peer\_\_cost\_ratio\_work\_type\_median} & $x_i / \text{Median}$ & Relative cost scaling \\
\texttt{peer\_\_cost\_work\_type\_robust\_z} & Robust Z on cost & Standardized cost anomaly \\
\texttt{peer\_\_unit\_cost\_ratio\_work\_type\_median} & $x_i^{unit} / \text{Median}^{unit}$ & Unit cost relative pricing \\
\texttt{peer\_\_unit\_cost\_work\_type\_robust\_z} & Robust Z on unit cost & Standardized unit overpricing \\
\texttt{peer\_\_cost\_deviation\_work\_type\_robust\_z} & Robust Z on cost deviation & Sanction-vs-estimate deviation \\
\texttt{peer\_\_payment\_velocity\_ratio\_work\_type\_median} & $v_i / \text{Median}(v)$ & Relative disbursement speed \\
\texttt{peer\_\_payment\_velocity\_work\_type\_robust\_z} & Robust Z on velocity & Standardized velocity anomaly \\
\texttt{peer\_\_payment\_concentration\_work\_type\_robust\_z} & Robust Z on concentration & Installment concentration anomaly \\
\texttt{peer\_\_contract\_value\_change\_work\_type\_robust\_z} & Robust Z on contract change & Post-award amendment anomaly \\
\texttt{peer\_\_cost\_category\_robust\_z} & Robust Z vs category & Category-level cost check \\
\bottomrule
\end{tabularx}
\end{table}


\vspace{1em}\hrule\vspace{1em}

\section{Topic 9 --- Normalization, Encoding, and Transformations}

\subsection{9.1 --- RobustScaler (why not StandardScaler)}

After 40 features are assembled (29 base + 11 peer), they are scaled with sklearn's \texttt{RobustScaler}:

\[ \tilde{x}_j = \frac{x_j - Q_2(x_j)}{Q_3(x_j) - Q_1(x_j)} \]

Where $Q_1, Q_2, Q_3$ are the 25th, 50th (median), 75th percentile of the \textbf{training population} for feature $j$.

\textbf{Why RobustScaler, not StandardScaler?}

Public works financial data is heavily right-skewed. A single \INR{}50 Crore highway project dominates the standard deviation of a work type where typical projects cost \INR{}5–15 Lakhs. StandardScaler's $\sigma$ would be completely distorted. RobustScaler uses IQR, which is unaffected by extreme values.

\subsection{9.2 --- Isolation Forest Score Normalization}

The Isolation Forest's raw output is a negative score (more negative = more anomalous). SETU converts it to 0–100:

\textbf{Step 1:} Negate:
\[ \text{raw\_anomaly}_i = -\text{decision\_function}(\mathbf{x}_i) \]

\textbf{Step 2:} Clip-normalize using population P1/P99:
\[ \text{anomaly\_score}_i = \text{clip}\!\left(\frac{\text{raw\_anomaly}_i - P_1}{P_{99} - P_1} \times 100,\ 0.0,\ 100.0\right) \]

$P_1$ and $P_{99}$ are computed on the training data and stored in the fitted model object:

\begin{lstlisting}[language=Python]
# isolation_forest_model.py lines 39–55
raw_anomaly = -self.model_.decision_function(X)
self.raw_score_p1_  = float(np.percentile(raw_anomaly, 1.0))
self.raw_score_p99_ = float(np.percentile(raw_anomaly, 99.0))
# At inference:
norm_score = (raw_anomaly - self.raw_score_p1_) / (self.raw_score_p99_ - self.raw_score_p1_) * 100.0
return np.clip(norm_score, 0.0, 100.0)
\end{lstlisting}


\subsection{9.3 --- Percentile Rank Computation}

Each project also receives a \textbf{population percentile rank}:

\[ \text{percentile}_i = \frac{\text{rank}(s_i)}{N} \times 100 \]

Using \texttt{scipy.stats.rankdata(method="average")} to handle ties. Percentile ≥ 95 $\rightarrow$ flagged anomaly.

\subsection{9.4 --- Temporal Transformations}

Raw timestamp columns are \textbf{never passed raw} into ML models:

\begin{table}[htbp]
\centering
\small
\begin{tabularx}{\textwidth}{l l Y}
\toprule
\textbf{Raw Timestamp} & \textbf{Derived Feature} & \textbf{Formula} \\
\midrule
\texttt{recommendation\_date} & \texttt{days\_since\_recommended} & $t_\text{today} - t_\text{rec}$ \\
\texttt{payment\_date\_dt} & \texttt{payment\_span\_days} & $t_\text{last\_payment} - t_\text{first\_payment}$ \\
\texttt{payment\_date\_dt} & \texttt{average\_days\_between\_payments} & Mean of $\Delta t$ between consecutive payments \\
\texttt{record\_date\_dt} & \texttt{physical\_progress\_slope} & $\Delta\text{progress} / \Delta\text{days}$ \\
\bottomrule
\end{tabularx}
\end{table}


\subsection{9.5 --- Missing Value Imputation}

\begin{lstlisting}[language=Python]
# preprocessor.py lines 191–193 — column median imputation
for col in combined_df.columns:
    val = combined_df[col].median()
    self.imputer_values_[col] = float(val) if not np.isnan(val) else 0.0
\end{lstlisting}


Median imputation (not mean) for the same reason as RobustScaler --- resistant to extreme financial outliers.

\subsection{9.6 --- Hard Negative Score Calibration (RAG Layer)}

After ML scoring, the RAG layer applies post-processing calibration:

\begin{lstlisting}[language=Python]
# hard_negative_evaluator.py lines 84–88
if portfolio_verdict == "HARD_NEGATIVE_MITIGATED":
    calibrated_score = max(20.0, raw_avg_risk * 0.55)
elif portfolio_verdict == "MIXED_EXPOSURE":
    calibrated_score = max(35.0, raw_avg_risk * 0.80)
# CONFIRMED_ANOMALY: no calibration
\end{lstlisting}


\textbf{HARD\_NEGATIVE\_MITIGATED} (×0.55): ≥60\% of flagged works have verified mitigations: flood zone + monsoon moratorium, statutory public agency designation, election MCC blackout, SoR-matching unit costs.

\textbf{MIXED\_EXPOSURE} (×0.80): Some works mitigated, some not --- partial adjustment.

\textbf{CONFIRMED\_ANOMALY}: Raw score stands unchanged.

\vspace{1em}\hrule\vspace{1em}

\section{Topic 10 --- Complete Data Lineage}

\subsection{Source $\rightarrow$ Model Input Trace}

\begin{tcolorbox}[colback=gray!5!white,colframe=gray!60!black,arc=1mm,boxrule=0.6pt,left=6pt,right=6pt,top=6pt,bottom=6pt]
\begin{lstlisting}[basicstyle=\ttfamily\scriptsize,breaklines=true,columns=flexible]
SOURCE                        TRANSFORMATION              MODEL INPUT
──────                        ──────────────              ───────────

01_projects.csv               prepare_project_base()      project__work_type
  (5,000 rows) ──────────►    prefix "project__"    ──►   project__category, project__status
                              strip: work_name,title       project__start_date ... (18 features)

02_financials.csv             prepare_prefixed_table()    financial__estimated_cost
  (5,000 rows) ──────────►    prefix "financial__"  ──►   financial__cost_per_unit
                                                          financial__cost_deviation
                                                          financial__sor_deviation ... (28 features)

03_payments.csv               aggregate_payments()        payment__payment_count
  (13,866 rows) ─────────►    → 5,000 rows                payment__payment_velocity_mean
                              prefix "payment__"    ──►   payment__unverified_payment_rate
                                                          payment__round_number_payment_rate
                                                          payment__payment_concentration_max
                                                          ... (29 features)

04_progress.csv               aggregate_progress()        progress__financial_minus_physical_progress
  (N rows) ──────────────►    → 5,000 rows                progress__physical_progress_slope
                              prefix "progress__"   ──►   progress__geotag_available_rate
                                                          progress__measurement_book_verified_rate
                                                          ... (26 features)

07_contractors.csv            extract_contractor_features() contractor__blacklisted_flag
  (400 rows, N:1 join) ──►    strip: contractor_name  ──►  contractor__capacity_utilization
                              prefix "contractor__"        contractor__shared_director_flag
                                                          ... (31 features)

09_geography.csv              extract_geography_features() geo__infrastructure_gap_index
  (619 rows, N:1 join) ──►    strip: district_name   ──►   geo__poverty_index
                              prefix "geo__"               geo__terrain_type ... (8 features)

12_labels.csv                 STRICT QUARANTINE            NEVER ENTERS MODEL
  (5,000 rows) ──────────►    → project_labels.csv  ──►   Used ONLY in post-hoc evaluator
\end{lstlisting}
\end{tcolorbox}


\subsection{Lineage --- Master Features to Model Input}

\begin{tcolorbox}[colback=gray!5!white,colframe=gray!60!black,arc=1mm,boxrule=0.6pt,left=6pt,right=6pt,top=6pt,bottom=6pt]
\begin{lstlisting}[basicstyle=\ttfamily\scriptsize,breaklines=true,columns=flexible]
project_master_features.csv (5,000 × 239)
    │
    └── FinancialPreprocessor.fit()
            ├── Select 29 base early-warning financial features
            ├── _compute_peer_stats()     → peer_stats_[work_type][col] = {median, MAD, count}
            ├── _engineer_peer_features() → 11 peer__ columns appended
            ├── Impute NaNs with column medians
            └── Fit RobustScaler on 40-column matrix
                    │
                    └── FinancialPreprocessor.transform()
                            → X_scaled (5,000 × 40 ndarray)
                                    │
                                    └── FinancialIsolationForestModel.fit(X_scaled)
                                            └── → financial_anomaly_scores.csv (5,000 × 17)
\end{lstlisting}
\end{tcolorbox}


\subsection{Lineage --- Domain Scores to API Response}

\begin{tcolorbox}[colback=gray!5!white,colframe=gray!60!black,arc=1mm,boxrule=0.6pt,left=6pt,right=6pt,top=6pt,bottom=6pt]
\begin{lstlisting}[basicstyle=\ttfamily\scriptsize,breaklines=true,columns=flexible]
financial_anomaly_scores.csv  (5,000 × 17)
geospatial_anomaly_scores.csv
procurement_anomaly_scores.csv       ┐
contractor_anomaly_scores.csv        ├── RiskFusionEngine.load_and_fuse_all_models()
payment_anomaly_scores.csv           │       → fuse_project() per row
progress_anomaly_scores.csv          │       → fused_risk_intelligence.csv (5,000 × ~25)
graph_anomaly_scores.csv             ┘               │
supervised_fraud_scores.csv                          └── FastAPI /risk-intelligence/projects
                                                              → JSON response to Next.js
\end{lstlisting}
\end{tcolorbox}


\subsection{Files Responsible for Each Transition}

\begin{table}[htbp]
\centering
\small
\begin{tabularx}{\textwidth}{l l Y}
\toprule
\textbf{Transition} & \textbf{File} & \textbf{Function} \\
\midrule
Raw CSVs $\rightarrow$ Validated tables & \texttt{validators.py} & \texttt{validate\_all()} \\
Payments (13,866) $\rightarrow$ Aggregated & \texttt{aggregators.py} & \texttt{aggregate\_payments()} \\
Progress $\rightarrow$ Aggregated & \texttt{aggregators.py} & \texttt{aggregate\_progress()} \\
12 tables $\rightarrow$ Master features (5,000×239) & \texttt{feature\_builder.py} & \texttt{build\_project\_master\_features()} \\
Master $\rightarrow$ Scaled financial input (5,000×40) & \texttt{preprocessor.py} & \texttt{fit\_transform()} \\
Scaled input $\rightarrow$ Anomaly scores & \texttt{isolation\_forest\_model.py} & \texttt{score()} \\
Anomaly scores $\rightarrow$ Reason traces & \texttt{isolation\_forest\_model.py} & \texttt{generate\_reason\_traces()} \\
8 domain scores $\rightarrow$ Fused risk & \texttt{risk\_fusion\_engine.py} & \texttt{load\_and\_fuse\_all\_models()} \\
RAG context $\rightarrow$ Hard negative calibration & \texttt{hard\_negative\_evaluator.py} & \texttt{evaluate\_portfolio()} \\
Fused CSV $\rightarrow$ API JSON & \texttt{risk\_intelligence.py} & \texttt{list\_fused\_projects()} \\
\bottomrule
\end{tabularx}
\end{table}


\vspace{1em}\hrule\vspace{1em}

\section{End-to-End Numerical Example --- Payment Aggregation}

\textbf{Scenario:} Project P-0042 has 4 payment installments:

\begin{table}[htbp]
\centering
\small
\begin{tabularx}{\textwidth}{c l c l Y}
\toprule
\textbf{Payment} & \textbf{Amount (\INR{})} & \textbf{Date} & \textbf{MB Verified} & \textbf{Round Number} \\
\midrule
1 & 2,00,000 & 2023-04-01 & Yes & No \\
2 & 1,50,000 & 2023-05-15 & Yes & No \\
3 & 3,00,000 & 2023-06-10 & No & No \\
4 & 14,90,000 & 2023-06-12 & No & Yes \\
\bottomrule
\end{tabularx}
\end{table}


\textbf{Aggregated values:}

\[ \text{pay\_count} = 4 \]
\[ \text{pay\_avg} = \frac{2{,}00{,}000 + 1{,}50{,}000 + 3{,}00{,}000 + 14{,}90{,}000}{4} = 5{,}35{,}000 \]
\[ \text{pay\_conc\_max} = \frac{14{,}90{,}000}{21{,}40{,}000} = 0.696 \]
\[ \text{pay\_span\_days} = (2023\text{-}06\text{-}12) - (2023\text{-}04\text{-}01) = 72 \text{ days} \]
\[ \text{pay\_velocity} = \frac{21{,}40{,}000}{72} = 29{,}722 \text{ \text{Rs.\,}/day} \]
\[ \text{unverified\_rate} = \frac{2}{4} = 0.50 \quad (50\%\ \text{unverified — critical flag}) \]
\[ \text{round\_rate} = \frac{1}{4} = 0.25 \quad (\text{\text{Rs.\,}14.9L just below \text{Rs.\,}15L threshold — structuring indicator}) \]

\textbf{After RobustScaler} (assuming population IQR for payment concentration = 0.20, median = 0.35):
\[ \tilde{x}_\text{conc} = \frac{0.696 - 0.35}{0.20} = 1.73 \]

This scaled value 1.73 enters the Isolation Forest as one of 40 input dimensions. Combined with other suspicious signals (unverified rate, round-number rate), the Isolation Forest scores this project near the 97th percentile $\rightarrow$ flagged as anomaly.

\vspace{1em}\hrule\vspace{1em}

\section{Important Equations --- Report 2 Summary}

\begin{table}[htbp]
\centering
\small
\begin{tabularx}{\textwidth}{l Y}
\toprule
\textbf{Equation} & \textbf{Source} \\
\midrule
$Z_i^\text{robust} = \frac{x_i - \text{Median}(\text{peer})}{1.4826 \times \text{MAD}(\text{peer}) + \varepsilon}$ & \texttt{preprocessor.py} L133 \\
$\text{pay\_velocity} = \frac{\sum a_j}{\max(1, t_\text{last} - t_\text{first})}$ & \texttt{aggregators.py} L73 \\
$\text{phys\_slope} = \frac{\Delta\text{physical\_progress}}{\Delta\text{days}}$ & \texttt{aggregators.py} L202 \\
$\tilde{x}_j = \frac{x_j - Q_2}{Q_3 - Q_1}$ (RobustScaler) & \texttt{preprocessor.py} L196 \\
$\text{anomaly\_score} = \text{clip}\!\left(\frac{-\text{df}(x) - P_1}{P_{99}-P_1} \times 100,\ 0,\ 100\right)$ & \texttt{isolation\_forest\_model.py} L54 \\
$\text{calibrated} = \max(20.0,\ \text{raw} \times 0.55)$ if HARD\_NEGATIVE\_MITIGATED & \texttt{hard\_negative\_evaluator.py} L86 \\
\bottomrule
\end{tabularx}
\end{table}


\vspace{1em}\hrule\vspace{1em}

\section{Common Misconceptions --- Report 2}

\begin{table}[htbp]
\centering
\small
\begin{tabularx}{\textwidth}{l Y}
\toprule
\textbf{Misconception} & \textbf{Reality} \\
\midrule
"The 12 tables are directly joined into the model" & 03\_payments and 04\_progress MUST be aggregated first. Direct joins produce duplicate project rows. \\
"StandardScaler is fine for financial data" & Financial data is heavily skewed. A \INR{}50Cr highway distorts σ. RobustScaler uses IQR, which is resistant. \\
"A high absolute cost means high risk" & Only suspicious relative to peers of the same work type. A \INR{}2Cr rural bridge is normal; \INR{}2Cr for a handpump is extreme. Peer normalization resolves this. \\
"Temporal features can be strings" & All dates are converted to elapsed-day intervals before entering models. \\
"The peer group is the full population" & Peer groups are work-type specific (29 types), then category (8), then global. Comparing a road to a handpump would be meaningless. \\
"Labels are used for training the domain models" & Labels are quarantined and only used post-hoc in \texttt{FinancialModelEvaluator} for benchmarking. Domain models are unsupervised. \\
\bottomrule
\end{tabularx}
\end{table}


\vspace{1em}\hrule\vspace{1em}

\section{Key Takeaways --- Report 2}

\begin{enumerate}[leftmargin=1.5em, itemsep=2pt, topsep=3pt]
\item \textbf{The many-to-one problem} is the fundamental data engineering challenge. Payments and progress are transaction-grain; models are project-grain. Aggregation is mandatory and deterministic.
\item \textbf{Peer normalization is the most critical feature engineering step.} Without it, large legitimate infrastructure projects dominate anomaly rankings purely due to scale.
\item \textbf{RobustScaler + MAD} are used throughout (not StandardScaler + σ) because MPLADS financial data is heavily right-skewed.
\item \textbf{Label quarantine is enforced at multiple independent levels}: schemas, feature\_builder assertions, preprocessor \texttt{\_validate\_safety()} --- all independently guard against leakage.
\item \textbf{Every output value can be traced back} to a specific row in a specific source CSV through the documented transformation chain.
\end{enumerate}

\vspace{1em}\hrule\vspace{1em}

\section{Prerequisites for Report 3}

Report 3 covers the mathematical and ML core: how Isolation Forest works as an isolation tree algorithm, the exact path length score derivation, why contamination=0.05 sets the decision boundary, the XGBoost dual-head classifier, SHAP-based explainability, and a complete numerical walkthrough tracing one real project through every transformation. Prerequisites from Reports 1–2: 40-dimensional RobustScaled feature vectors, normalized anomaly scores [0,100], work-type peer groups, quarantined labels.

\vspace{1em}\hrule\vspace{1em}

\textit{Say \textbf{"Continue"} for Report 3 --- Mathematical \& ML Core (Isolation Forest internals, XGBoost, SHAP)}
\textit{Say \textbf{"Go deeper"} on any specific topic in Report 1 or 2}
\textit{Say \textbf{"Trace this"} + a feature name to trace its full journey from source CSV to model output}

\vspace{1em}\hrule\vspace{1em}

\textbf{SETU --- Smart Evidence-based Triangulation for Uncovering Anomalies in MPLADS}
\textit{Reports prepared from source-code analysis of \texttt{backend/app/ml/} and \texttt{backend/app/api/}.}
\textit{Developed for Smart India Hackathon 2026 --- MoSPI Problem Statement.}
\end{document}
